% Part: lambda-calculus
% Chapter: lambda-definability
% Section: arithmetical-functions

\documentclass[../../../include/open-logic-section]{subfiles}

\begin{document}

\olfileid{lam}{rep}{arf}
\olsection{Fungsi Aritmetika kang \usetoken{S}{lambda definable}}

\begin{prop}
  \ollabel{prop:succ-ld}
  Fungsi penerus~$\Succ$ iku
  !!{lambda definable}.
\end{prop}

\begin{proof}
Suku kang !!{lambda define}s
fungsi penerus yaiku
\begin{align*}
  \fn{Succ} & \ident \lambd[a][\lambd[fx][f ({a} f x)]].
  \intertext{Miturut konvensi kita, iki cekakan saka}
  \fn{Succ} & \ident \lambd[a][\lambd[f][\lambd[x][(f (({a} f) x))]]].
  \intertext{$\fn{Succ}$ yaiku fungsi kang nampa
    sawijining bilangan~$a$ minangka argumen,
    banjur ngasilake fungsi liya,
    $\lambd[fx][f ({a} f x)]$.
    Fungsi iki dhewe dudu numeral Church.
    Nanging yen argumen $a$ numeral Church,
    asile direduksi dadi numeral Church.
    Gatekna:}
   (\lambd[a][\lambd[fx][f ({a} f x)]])\,\num{n} & \redone
   \lambd[fx][f ({\num{n}} f x)].
   \intertext{Suku sisipan $\num{n}fx$ iku redex,
     amarga $\num{n}$ yaiku $\lambd[fx][f^nx]$.
     Mula $\num{n}fx \redone f^nx$,
     saéngga kanggo suku sakabèhé kita oleh}
   \fn{Succ}\, \num n & \red \lambd[fx][f(f^n(x))],
\end{align*}
yaiku $\num{n+1}$.
\end{proof}

\begin{ex}
  Gatekna apa kang kelakon yen
  $\fn{Succ}$ ditrapake
  marang~$\num{0}$, yaiku
  $\lambd[fx][x]$.
  Kita tulis suku-sukune
  kanthi pepak:
  \begin{align*}
    \fn{Succ}\, \num{0} & \ident (\lambd[a][\lambd[f][\lambd[x][(f (({a} f) x))]]])(\lambd[f][\lambd[x][x]])\\
    & \redone \lambd[f][\lambd[x][(f (({(\lambd[f][\lambd[x][x]])} f) x))]] \\
    & \redone \lambd[f][\lambd[x][(f ({(\lambd[x][x])} x))]] \\
    & \redone \lambd[f][\lambd[x][(f x)]] \ident \num{1}\\
  \end{align*}
\end{ex}

\begin{prob}
  Suku
  \begin{align*}
    \fn{Succ}' & \ident \lambd[{n}][\lambd[fx][{n} f (f x)]]
  \end{align*}
  !!{lambda define}s
  fungsi penerus.
  Terangna alesane.
\end{prob}

\begin{prop}
  \ollabel{prop:add-ld}
  Fungsi panjumlahan~$\Add$
  iku !!{lambda definable}.
\end{prop}

\begin{proof}
Panjumlahan !!{lambda define}d
dening suku-suku iki:
\begin{align*}
  \fn{Add} & \ident \lambd[{a}{b}][\lambd[fx][{a} f ({b} f x)]]
\intertext{utawa, minangka cara liya,}
  \fn{Add}' & \ident \lambd[{a}{b}][{a}\, \fn{Succ} \, {b}].
\intertext{Panjumlahan kapisan lumaku mangkene:
  $\fn{Add}$ nampa dhisik rong bilangan,
  ${a}$ lan ${b}$. Asile sawijining fungsi
  kang nampa $f$ lan $x$ banjur
  ngasilake $af(bfx)$. Yen $a$ lan $b$
  numeral Church $\num{n}$ lan
  $\num{m}$, asil iki direduksi dadi
  $f^{n+m}(x)$, kang padha karo
  $f^{n}(f^{m}(x))$. Rinciane:}
  (\lambd[{a}{b}][\lambd[fx][{a} f ({b} f x)]])\num n\,\num m & \redone
  \lambd[fx][\num{n}\, f (\num {m}\, f x)] \\
  & \redone \lambd[fx][\num{n}\, f (f^m x)] \\
  & \redone \lambd[fx][f^n (f^m x)] \ident \num {n+m}.
\intertext{Perwakilan panjumlahan kapindho,
  $\fn{Add'}$, lumaku kanthi cara liya:
  yen ditrapake marang rong numeral Church,
  $\num{n}$ lan $\num{m}$,}
\fn{Add}' \num n \,\num m
& \redone \num{n}\, \fn{Succ}\, \num{m}.
\intertext{Nanging $\num{n} f x$ mesthi
  direduksi dadi $f^n(x)$. Mula,}
  \num{n}\,  \fn{Succ}\, \num{m} & \red \fn{Succ}^n(\num{m}).
\end{align*}
Amarga $\fn{Succ}$
!!{lambda define}s fungsi
penerus, lan fungsi penerus
kang ditrapake ping $n$
marang~$m$ ngasilake
$n+m$, asil iki banjur
direduksi dadi~$\num{n+m}$.
\end{proof}

\begin{prop}
  \ollabel{prop:mult-ld}
  Perkalian iku
  !!{lambda definable}
  dening suku iki:
  \[
  \fn{Mult} \ident \lambd[ab][\lambd[fx][a (b f) x]]
  \]
\end{prop}

\begin{proof}
  Kanggo ndeleng carane kerjane,
  upamane $\fn{Mult}$
  ditrapake marang numeral Church
  $\num{n}$ lan $\num{m}$:
  $\fn{Mult} \, \num{n} \, \num{m}$
  direduksi dadi
  $\lambd[fx][\num{n}(\num{m}\, f)x]$.
  Suku $\num{m} f$ netepake
  fungsi kang ngetrapake $f$
  marang argumene ping $m$.
  Mula $\num{n} (\num{m} f) x$
  ngetrapake fungsi kang
  ngetrapake $f$ ping $m$
  iku dhewe ping $n$
  marang~$x$.
  Kanthi tembung liya,
  $f$ ditrapake marang
  $x$ ping $n\cdot m$.
  Suku normal asile
  ora liya numeral Church
  $\num{nm}$.
\end{proof}

\begin{editorial}
Suku iki sejatine bisa
disederhanakake maneh
kanthi reduksi-$\eta$:
\[
  \fn{Mult} \ident \lambd[ab][\lambd[f][a (b f)]].
\]

Reduksi-$\eta$ kang
dibutuhake ing kene wis
diandharake ing bab sintaksis.
\end{editorial}

\begin{prob}
Perkalian bisa
!!{lambda define}d
dening suku
\[
  \fn{Mult}' \ident \lambd[ab][a (\fn{Add}\, b) \num{0}].
\]
Terangna kepriye
suku iki makarya.
\end{prob}

Definisi pemangkatan
minangka suku-$\lambd$
prasaja kanthi cara
kang nggumunake:
\[
  \fn{Exp} \ident \lambd[be][e b].
\]
Argumen kapisan $b$
yaiku basis, dene
argumen kapindho~$e$
yaiku eksponen.
Miturut intuisi,
$e f$ iku $f^e$
miturut pangodean
bilangan kita.
Yen iki angel dipahami,
pemangkatan uga bisa
ditetepake liwat
perkalian kang dibaleni:
\[
  \fn{Exp}' \ident \lambd[be][e (\fn{Mult}\, b) \num{1}].
\]

Fungsi sadurunge lan
pengurangan ing numeral
Church ora segampang
kang dikira: perlu
pangodean pasangan.

\end{document}
